
\section{本章内容}

本体定义了知识的"骨架"（TBox），知识图谱在其上填充大规模实例数据（ABox）。
本章讲述从企业数据源到可用知识图谱的工程全流程：
\begin{itemize}
\item 构建流水线：数据接入 \(\rightarrow\) 实体/关系抽取 \(\rightarrow\) 知识融合 \(\rightarrow\) 入库
\item 实体对齐与消歧：同一设备的多种叫法如何归一
\item 数据质量：用 SHACL 形状约束做自动化校验
\item 存储与查询：三元组库与属性图数据库的选型
\end{itemize}

\section{文件说明}

\begin{center}
\begin{tabularx}{\linewidth}{@{}XX@{}}
\toprule
\textbf{文件} & \textbf{内容} \\
\midrule
\texttt{kg-construction-pipeline.txt} & 从台账/工艺文件/维修日志到知识图谱的流水线 \\
\texttt{entity-resolution.txt} & 实体对齐与消歧：规则、相似度、冲突消解 \\
\texttt{kg-quality-shacl.ttl} & SHACL 数据质量校验形状（可直接用 pySHACL 运行） \\
\texttt{kg-storage-query.txt} & 三元组库 vs 属性图选型；SPARQL 与 Cypher 对照 \\
\bottomrule
\end{tabularx}
\end{center}

\section{本体与知识图谱的关系}

\begin{center}
\begin{tabularx}{\linewidth}{@{}XXll@{}}
\toprule
\textbf{层次} & \textbf{内容} & \textbf{规模} & \textbf{维护者} \\
\midrule
本体 (TBox) & 类、属性、公理 & 数十\textasciitilde{}数百概念 & 本体工程师+专家 \\
知识图谱 (TBox+ABox) & 本体 + 海量实例与关系 & 百万\textasciitilde{}十亿三元组 & 自动化管道 \\
\bottomrule
\end{tabularx}
\end{center}

\section{技术栈}

\begin{center}
\begin{tabular}{@{}ll@{}}
\toprule
\textbf{环节} & \textbf{代表工具} \\
\midrule
实体/关系抽取 & 规则模板、NER模型、LLM抽取（需校验，见第8章） \\
质量校验 & SHACL (pySHACL / TopBraid)、ShEx \\
三元组存储 & Jena TDB2 + Fuseki、GraphDB、Virtuoso \\
属性图存储 & Neo4j、NebulaGraph \\
\bottomrule
\end{tabular}
\end{center}
